
\def\PUDcodigo{ECA220}

\def\PUDcodacad{04507.26}

\def\PUDdisciplina{Sistemas Lineares}

\def\PUDsemestre{S5}

\def\PUDhoras{80}

%NOVOS ---------------------------------------------------
\def\PUDhorasTeoria{72}
\def\PUDhorasPratica{8}

\def\PUDinterECA{ \hyperref[PUD:ACE002]{Álgebra Linear (04507.1)} 

\hyperref[PUD:ECA210]{Circuitos Elétricos I (04507.20)}

\hyperref[PUD:ECA233]{Controle I (04507.31)}
}

\def\PUDatuaisECA{- Apresentação de softwares baseados em matrizes, como Matlab, SciLab ou Octave}

\def\PUDrecursos{Projetor multimídia; Quadro branco e pincel.}

%----------------------------------------------------------

\def\PUDcreditos{4}

\def\PUDprerequisito{ \hyperref[PUD:ACE002]{ACE002} }

\def\PUDpreacad{ \hyperref[PUD:ACE002]{Álgebra Linear (04507.1)} }

\def\PUDobjetivos{Compreender a funções que descrevem 
sinais e sistemas lineares.  Analisar sistemas dinâmicos
 contínuos e discretos.}

\def\PUDementa{Introdução aos sistemas LTI;  Representação de sinais periódicos em séries;  Transformadas de tempo contínuo e de tempo discreto;  Sistemas lineares com realimentação.}

\def\PUDconteudo{ & Introdução aos sistemas de tempo contínuo e discreto, transformação de variáveis independente.  Funções de impulso e degrau unitário. 
Propriedades básicas de sistemas.  
\\ 
 & Sistemas lineares invariantes no tempo.  Soma de convolução.  Integral de convolução.  Propriedades dos sistemas LIT.  Sistemas LIT causais
descritos por equações diferenciais e de diferença.  Funções de singularidade. 
\\ 
 & Representação de sinais periódicos em série de Fourier.  Resposta dos sistemas LIT.  Representações de tempo contínuo e discreto. 
Filtragem e filtros de tempo contínuo e de tempo discreto. 
\\ 
 & Transformada de Fourier de tempo contínuo.  propriedades. 
\\ 
 & Transformada de Fourier de tempo discreto.  propriedades. 
\\ 
 & Caracterização no tempo e na frequência dos sinais e sistemas.  Representação magnitude e fase da transformada de Fourier e da 
resposta em frequência dos sistemas LIT.  Sistemas de primeira ordem de tempo contínuo e discreto. }

\def\PUDmetodo{Aulas expositórias que podem ser teóricas e/ou práticas, onde as práticas no laboratório serão marcadas no decorrer da disciplina.}

\def\PUDavalia{A avaliação será feita com aplicação de provas teóricas e/ou práticas, além da possibilidade de inclusão de trabalhos e seminários no decorrer da disciplina.}

\def\PUDbibbasica{ & \href{www.biblioteca.ifce.edu.br/index.asp?codigo_sophia=39861}{Oppenhein}, Alan. V.; Willsky, A. S.; Hamid, S.  Sinais e sistemas.  2º ed.  São Paulo, SP: Pearson Prentice Hall, 2010.  568p.  ISBN: 9788576055044.
\\ 
 &  \href{www.biblioteca.ifce.edu.br/index.asp?codigo_sophia=39837}{Lathi}, B. P.  Sinais e Sistemas Lineares.  2º ed.  Porto Alegre, RS: Bookman, 2007.  856p.  ISBN: 9788560031139.
\\ 
 %&  HSU, Hwei P.  Sinais e Sistemas - Coleção Shaum.  2a ed.  Porto Alegre(RS): Editora Bookman.  2012.  508p.  ISBN 9788577809387. }
 &  \href{www.biblioteca.ifce.edu.br/index.asp?codigo_sophia=40232}{Oppenheim}, Alan V.  Processamento em tempo discreto de sinais.  3º ed.  São Paulo, SP: Pearson Education do Brasil, 2012.  665p.  ISBN: 9788581431024. }

\def\PUDbibcomplementar{ &  \href{www.biblioteca.ifce.edu.br/index.asp?codigo_sophia=63268}{Roberts}, M. J.  Fundamentos em Sinais e Sistemas.  São Paulo, SP: McGraw-Hill, 2009.  764p.  ISBN: 9788577260386.
\\ 
 &  \href{www.biblioteca.ifce.edu.br/index.asp?codigo_sophia=62379}{Girod}, Bernd.  Sinais e Sistemas.  Rio de Janeiro, RJ: LTC, 2003.  340p.  ISBN: 8521613644.
\\ 
 %&  CHAPARRO, Luis.  Signals and System using MATLAB.  ACADEMIC PRESS.  2010.  650p.  ISBN : 9780123747167.
 &  \href{www.biblioteca.ifce.edu.br/index.asp?codigo_sophia=24005}{Oliveira}, André Schneider de.  Sistemas embarcados: hardware e firmware na prática.  São Paulo, SP: Érica, 2006.  316p.  ISBN: 8536501057.
\\
 & \href{www.biblioteca.ifce.edu.br/index.asp?codigo_sophia=24455}{Barroso}, Leonidas.  Cálculo numérico.  2º ed.  São Paulo, SP: Harbra, 1987.  367p.  ISBN: 8529400895.
\\
 & \href{www.biblioteca.ifce.edu.br/index.asp?codigo_sophia=62475}{Gilat}, Amos.  MATLAB com aplicações em engenharia.  4º ed.  Porto Alegre, RS: Bookman, 2012.  417p.  ISBN: 9788540701861. }

\def\PUDautor{Adriano Pereira}

\def\PUDdata{2013-04-29}

\def\PUDrevisao{3}

\def\PUDrevisor{Adriano Holanda Pereira}

\def\PUDdatarevisao{2019-05-15}

\PUDcorpo
